\documentclass[aspectratio=169,12pt]{beamer}
\usepackage{fontspec}
\usepackage[portuguese]{babel}
\usepackage{xcolor,listings,tikz,booktabs,amsmath,amssymb,graphicx,multicol,array,colortbl}
\usetikzlibrary{arrows.meta,positioning,shapes.geometric,calc}

\definecolor{navy}{HTML}{0D1B3E}
\definecolor{gold}{HTML}{F5A623}
\definecolor{qgreen}{HTML}{1E8B4C}
\definecolor{qblue}{HTML}{1A6EAE}
\definecolor{codebg}{HTML}{EEF3F8}

\usetheme{default}
\useinnertheme{circles}
\setbeamertemplate{navigation symbols}{}
\setbeamertemplate{footline}{%
  \hfill{\color{gray!60}\scriptsize\insertframenumber/\inserttotalframenumber}\hspace{.6em}\vspace{.4em}}

\setbeamercolor{frametitle}{bg=navy,fg=white}
\setbeamercolor{structure}{fg=navy}
\setbeamercolor{alerted text}{fg=gold}
\setbeamercolor{item}{fg=gold}
\setbeamercolor{subitem}{fg=navy}
\setbeamercolor{block title}{bg=qblue,fg=white}
\setbeamercolor{block body}{bg=qblue!8,fg=black}
\setbeamercolor{block title alerted}{bg=gold,fg=navy}
\setbeamercolor{block body alerted}{bg=gold!15,fg=navy}
\setbeamercolor{block title example}{bg=qgreen,fg=white}
\setbeamercolor{block body example}{bg=qgreen!10,fg=black}
\setbeamerfont{frametitle}{size=\large,series=\bfseries}
\setbeamerfont{block title}{series=\bfseries}

\setbeamertemplate{title page}{%
  \begin{tikzpicture}[remember picture,overlay]
    \fill[navy](current page.north west)rectangle(current page.south east);
    \fill[gold]([yshift=.7cm]current page.south west)rectangle(current page.south east);
    \node[anchor=south,yshift=.73cm,navy,font=\tiny\bfseries]at(current page.south)
      {INTENSIVÃO QUANT AI \textbullet{} IMPACT UFSCAR};
  \end{tikzpicture}
  \vspace{.65cm}
  \begin{center}
    {\color{gold}\scriptsize\bfseries INTENSIVÃO QUANT AI \textbullet{} IMPACT UFSCAR}\\[.4cm]
    {\color{white}\LARGE\bfseries\inserttitle}\\[.25cm]
    {\color{gold!85}\normalsize\insertsubtitle}\\[.5cm]
    {\color{white!70}\small\insertauthor}\\[.1cm]
    {\color{white!50}\footnotesize\insertdate}
  \end{center}}

\lstdefinestyle{python}{language=Python,
  basicstyle=\ttfamily\scriptsize\color{qblue},
  keywordstyle=\color{navy}\bfseries,
  stringstyle=\color{qgreen!80!black},
  commentstyle=\color{gray!70}\itshape,
  backgroundcolor=\color{codebg},
  frame=l,framerule=2pt,rulecolor=\color{gold},
  xleftmargin=1em,breaklines=true,showstringspaces=false,tabsize=4,
  extendedchars=false,inputencoding=ascii}
\lstset{style=python}

\author{Felipe Maldonado\\{\scriptsize VP Diretoria Quant~\textbullet{}~Impact UFSCAR}}
\date{Intensivão Quant AI 2026}
\title{Aula 02 --- Dados}
\subtitle{Coleta, Limpeza e Preparação de Dados Financeiros}

\begin{document}

\begin{frame}[plain]
\titlepage
\end{frame}

\begin{frame}[fragile]{Nesta Aula}
\begin{columns}[T]
  \column{0.5\textwidth}
\begin{block}{Teoria (20 min)}
\begin{itemize}
  \item O universo IBOVESPA: composição e rebalanceamento
  \item Pipeline de dados: da fonte ao parquet
  \item Tratamento de outliers e dados faltantes
  \item Retornos compostos vs simples
\end{itemize}
\end{block}
  \column{0.47\textwidth}
\begin{block}{Código (40 min)}
\begin{itemize}
  \item \texttt{baixar\_dados()} — yfinance $\to$ DataFrame
  \item \texttt{limpar\_precos()} — outliers, forward-fill
  \item \texttt{calcular\_retornos\_mensais()} — resample + log-ret
  \item Salvar 3 parquets base
  \item Verificação de sanidade dos dados
\end{itemize}
\end{block}
\end{columns}
\end{frame}

\begin{frame}[fragile]{O Universo IBOVESPA}
\begin{columns}[T]
  \column{0.5\textwidth}
\begin{itemize}
  \item Índice de retorno total da B3: ~90 ativos
  \item Critério de inclusão: liquidez (VDLA $\geq$ 0{,}1\% do total)
  \item Rebalanceamento: janeiro, maio e setembro
  \item Composição histórica obtida via yfinance
  \item Período de análise: 2012--2024 (12 anos)
\end{itemize}
  \column{0.47\textwidth}
\begin{block}{Por que o IBOVESPA?}
\begin{itemize}
  \item Universo amplo o suficiente para cross-section
  \item Dados históricos de qualidade disponíveis
  \item Relevante para o Desafio Itaú
  \item Benchmark natural para estratégias Long Only
\end{itemize}
\end{block}\vspace{.4em}
\begin{alertblock}{Atenção: Survivorship Bias}
\begin{itemize}
  \item yfinance retorna composição atual + histórico de preços
  \item Ativos que saíram do índice ainda têm dados históricos
  \item Bias reduzido, mas não eliminado — limitação declarada
\end{itemize}
\end{alertblock}
\end{columns}
\end{frame}

\begin{frame}[fragile]{Pipeline de Dados}

\begin{center}
\begin{tikzpicture}[node distance=.6cm and .4cm,
  box/.style={draw=navy,fill=navy!8,rounded corners=3pt,
              text width=2.2cm,align=center,font=\scriptsize\bfseries,inner sep=5pt},
  arr/.style={-Stealth,navy,thick}]
  \node[box](a){yfinance\\{\tiny download}};
  \node[box,right=of a](b){precos\_ibov\\{\tiny parquet}};
  \node[box,right=of b](c){limpar\\{\tiny outliers}};
  \node[box,right=of c](d){ret diários\\{\tiny parquet}};
  \node[box,right=of d](e){resample\\{\tiny mensal}};
  \node[box,right=of e](f){ret mensais\\{\tiny parquet}};
  \draw[arr](a)--(b); \draw[arr](b)--(c); \draw[arr](c)--(d);
  \draw[arr](d)--(e); \draw[arr](e)--(f);
\end{tikzpicture}
\end{center}
\vspace{.3em}
\begin{columns}[T]
  \column{0.48\textwidth}
  \begin{block}{Outlier Detection}
    Retorno diário $|r_t| > 3\sigma$ rolling 60d $\Rightarrow$ substituído por NaN \\
    Forward-fill limitado a 5 dias consecutivos
  \end{block}
  \column{0.48\textwidth}
  \begin{block}{Retorno Composto Mensal}
    $r_\text{mês} = \prod_{t \in \text{mês}}(1+r_t) - 1$ \\
    Equivalente a \texttt{resample('ME').apply(retorno\_composto)}
  \end{block}
\end{columns}
\end{frame}

\begin{frame}[fragile]{O que Vamos Construir}
\begin{columns}[T]
  \column{0.52\textwidth}
\begin{block}{Funções a implementar}
\begin{itemize}
  \item \texttt{baixar\_dados(tickers, start, end)} $\to$ DataFrame de preços
  \item \texttt{limpar\_precos(df)} $\to$ DataFrame limpo, outliers removidos
  \item \texttt{calcular\_retornos\_mensais(precos)} $\to$ DataFrame mensal
  \item \texttt{verificar\_sanidade(df)} $\to$ estatísticas de qualidade
\end{itemize}
\end{block}
  \column{0.45\textwidth}
\begin{block}{Entradas (parquets)}
\begin{itemize}
  \item Nenhum (download direto via yfinance)
\end{itemize}
\end{block}
\vspace{.4em}
\begin{exampleblock}{Saídas (parquets)}
\begin{itemize}
  \item \texttt{precos\_ibov.parquet}
  \item \texttt{retornos\_diarios\_limpo.parquet}
  \item \texttt{retornos\_mensais\_limpo.parquet}
\end{itemize}
\end{exampleblock}
\end{columns}
\end{frame}

\begin{frame}[fragile]{Referências}
\begin{multicols}{2}
\tiny
\begin{itemize}
  \item \textbf{Cont, R.} (2001). Empirical properties of asset returns: stylized facts and statistical issues. \textit{Quantitative Finance}, 1(2), 223--236.
  \item \textbf{Harvey, C. et al.} (2016). ... and the cross-section of expected returns. \textit{RFS}, 29(1), 5--68.
  \item \textbf{yfinance} (2019). Yahoo Finance Python wrapper. \textit{GitHub: ranaroussi/yfinance}.
  \item \textbf{Mussa, A. et al.} (2012). Hipótese de mercados eficientes e finanças comportamentais. \textit{REGE}, 19(2).
\end{itemize}
\end{multicols}
\end{frame}


\end{document}
